Trobem una aplicació de la inducció electromagnètica en els aparells de soldadura elèctrica. En un d'aquests aparells desmuntat veiem dues bobines com les d'un transformador.

La bobina primària té 1\,000 espires i la secundària en té 20. En la bobina secundària, feta d'un fil molt més gruixut, és on va connectat l'elèctrode per a fer la soldadura.

Sabem, per les especificacions tècniques impreses en la màquina, que pel circuit secundari circula una intensitat de corrent de 100~A. Determineu:

\begin{apartat}{1}
La tensió del circuit secundari quan es connecta la màquina, és a dir, quan es connecta el circuit primari a una tensió alterna de 220~V.
\end{apartat}

\begin{apartat}{1}
La intensitat que circula pel circuit primari i la potència consumida per la màquina.
\end{apartat}

\nota{Negligiu qualsevol tipus de dissipació d'energia.}
